\documentclass[aspectratio=169, 10pt]{beamer}

\usepackage[utf8]{inputenc}
\usepackage[T5]{fontenc}
\usepackage[vietnamese]{babel}
\usepackage{times}
\usepackage{booktabs}
\usepackage{xcolor}
\usepackage{tcolorbox}
\tcbuselibrary{skins}

\definecolor{primaryDark}{RGB}{0, 51, 102}
\definecolor{primaryLight}{RGB}{0, 85, 150}
\definecolor{accentTeal}{RGB}{0, 100, 100}
\definecolor{alertRed}{RGB}{220, 50, 50}
\definecolor{doneGreen}{RGB}{0, 130, 60}
\definecolor{textBlack}{RGB}{30, 30, 30}
\definecolor{lightGrayBg}{RGB}{245, 245, 245}

\usetheme{Madrid}
\usecolortheme{default}

\setbeamercolor{structure}{fg=primaryDark}
\setbeamercolor{palette primary}{bg=primaryDark, fg=white}
\setbeamercolor{palette secondary}{bg=primaryLight, fg=white}
\setbeamercolor{palette tertiary}{bg=accentTeal, fg=white}
\setbeamercolor{frametitle}{bg=primaryDark, fg=white}
\setbeamercolor{title}{bg=primaryDark, fg=white}
\setbeamercolor{block title}{bg=primaryDark, fg=white}
\setbeamercolor{block body}{bg=lightGrayBg, fg=textBlack}
\setbeamercolor{block title alerted}{bg=alertRed, fg=white}
\setbeamercolor{block title example}{bg=doneGreen, fg=white}
\setbeamercolor{alerted text}{fg=alertRed}

\setbeamertemplate{navigation symbols}{}
\setbeamertemplate{blocks}[default]
\setbeamertemplate{itemize item}{\color{accentTeal}$\blacksquare$}
\setbeamertemplate{itemize subitem}{\color{accentTeal}$\bullet$}
\setbeamertemplate{enumerate item}{\color{primaryDark}\textbf{\insertenumlabel.}}

\title[SDN \& SD-WAN Campus]{\texorpdfstring{XÂY DỰNG MẠNG CAMPUS NETWORK\\SỬ DỤNG SDN VÀ SD-WAN}{XÂY DỰNG MẠNG CAMPUS NETWORK SỬ DỤNG SDN VÀ SD-WAN}}
\subtitle{Báo cáo tiến độ hằng tuần -- Dự án Công nghệ Thông tin}
\author[Trần Hữu Nhân -- Nguyễn Nhật Hào]{Trần Hữu Nhân -- 52300235\\Nguyễn Nhật Hào -- 52300198}
\institute[TDTU]{Khoa Công nghệ Thông tin -- Trường Đại học Tôn Đức Thắng}
\date{\today}

\begin{document}

% ============================================================
% TITLE
% ============================================================
\begin{frame}[plain]
    \titlepage
\end{frame}

\begin{frame}{Nội dung báo cáo}
    \tableofcontents
\end{frame}

% ============================================================
% OBJECTIVE
% ============================================================
\section{Mục tiêu}

\begin{frame}{Mục tiêu đề tài}
    \begin{block}{Mục tiêu tổng quát}
        Xây dựng mô hình mạng Campus đa cơ sở cho trường đại học trên EVE-NG, kết hợp
        \textbf{SDN} (điều khiển tập trung mạng LAN campus chính) và
        \textbf{SD-WAN} (kết nối thông minh giữa các cơ sở qua Internet và MPLS).
    \end{block}
    \vspace{0.3cm}
    \begin{itemize}
        \item Thiết kế kiến trúc phân cấp Core -- Distribution -- Access, dự phòng VRRP, OSPF Area 0.
        \item Quy hoạch VLAN, IP, DHCP cho Site 100 (Campus chính) và 3 chi nhánh
              Cần Thơ (200), Đà Nẵng (300), Nha Trang (400).
        \item Quản lý tập trung các switch OVS bằng Ryu Controller qua OpenFlow 1.3.
        \item Dựng fabric Cisco SD-WAN (vManage/vSmart/vBond -- Site 900) với 8 vEdge,
              underlay BGP 2 nhà cung cấp (Internet AS 64511, MPLS AS 64512).
        \item Bảo mật và dịch vụ: FW-ASAv HA, DMZ, Server Farm, Syslog, giám sát NOC.
    \end{itemize}
\end{frame}

\begin{frame}{Mục tiêu của tuần}
    \begin{enumerate}
        \item Ổn định fabric SD-WAN: đủ 8/8 vEdge \emph{reachable}, BFD/tunnel full-mesh lên đủ 2 màu.
        \item Khắc phục đường dữ liệu SDN campus để gói DHCP đi xuyên các switch OVS.
        \item Đưa giao diện quản trị vManage ra mạng LAN thật để thao tác từ laptop.
        \item Phân tích, giải thích số liệu tunnel và Transport Health trên vManage.
        \item Đồng bộ cấu hình lab (host 1 $\rightarrow$ repo $\rightarrow$ host 2) và cập nhật tài liệu.
    \end{enumerate}
\end{frame}

% ============================================================
% PROGRESS
% ============================================================
\section{Tiến độ đã thực hiện}

\begin{frame}{Tổng quan tiến độ theo hạng mục}
    \centering
    \small
    \begin{tabular}{@{}llc@{}}
        \toprule
        \textbf{Giai đoạn} & \textbf{Hạng mục} & \textbf{Trạng thái} \\
        \midrule
        Khảo sát \& thiết kế    & Lý thuyết SDN/SD-WAN, topo, bảng IP/VLAN/ASN   & \textcolor{doneGreen}{\textbf{Hoàn thành}} \\
        Hạ tầng campus L2/L3    & Core/Dist/Access, VRRP, OSPF, FW-ASAv HA, ASDM & \textcolor{doneGreen}{\textbf{Hoàn thành}} \\
        Chi nhánh 200/300/400   & Brand-FW, DHCP, SwitchBrand, VPC              & \textcolor{doneGreen}{\textbf{Hoàn thành}} \\
        SD-WAN                  & Controllers, PKI, 8 vEdge, BGP underlay        & \textcolor{doneGreen}{\textbf{Hoàn thành}} \\
        SDN campus chính        & Ryu + 6 OVS, app L2, NOC Dashboard             & \textcolor{orange}{\textbf{Đang làm}} \\
        DHCP campus chính       & Win Server 2012 R2, 4 scope, relay trên Core   & \textcolor{orange}{\textbf{Đang làm}} \\
        Dịch vụ \& kiểm thử     & Web/Mail, ma trận test, failover               & \textcolor{alertRed}{\textbf{Chưa làm}} \\
        \bottomrule
    \end{tabular}
\end{frame}

\begin{frame}{Tiến độ -- SD-WAN}
    \begin{columns}[T]
        \begin{column}{0.5\textwidth}
            \begin{exampleblock}{Đã hoàn thành}
                \begin{itemize}
                    \item 3 controller (vManage/vSmart/vBond) + CA gốc, quy trình CSR/ký/cài root-cert-chain.
                    \item Whitelist đủ 8 serial vEdge trên \textbf{cả 3 controller}.
                    \item 8/8 vEdge \emph{reachable}; 10 TLOC, 36 tunnel duy nhất (72 phiên BFD).
                    \item Underlay BGP: backbone eBGP 64511 $\leftrightarrow$ 64512 \emph{Established}, CE--PE trên 8 vEdge.
                \end{itemize}
            \end{exampleblock}
        \end{column}
        \begin{column}{0.5\textwidth}
            \begin{block}{Sự cố đã xử lý trong tuần}
                \begin{itemize}
                    \item vEdge S100 chỉ \emph{1/3 vSmart}: thêm route /24 tới controller qua ngả
                          biz-internet (ECMP) $\Rightarrow$ đủ 3/3 control.
                    \item vEdge S100 rơi fabric sau reboot: bổ sung lại whitelist valid-vedges.
                    \item Switch61 nối Cloud $\rightarrow$ LAN thật: truy cập vManage GUI từ laptop.
                    \item Viết tài liệu giải thích Tunnel \& Transport Health.
                \end{itemize}
            \end{block}
        \end{column}
    \end{columns}
\end{frame}

\begin{frame}{Tiến độ -- SDN campus chính \& hạ tầng}
    \begin{columns}[T]
        \begin{column}{0.5\textwidth}
            \begin{block}{SDN (Ryu + Open vSwitch)}
                \begin{itemize}
                    \item Ryu chạy trên SDN\_CONTROLLER (10.1.99.10, port 6653/8080), điều khiển qua VLAN 99.
                    \item Sửa app \texttt{campus\_switch\_13.py}: table-miss
                          \texttt{CONTROLLER} $\rightarrow$ \texttt{NORMAL}, tự cài lại khi switch kết nối.
                    \item DHCP Discover đã đi xuyên campus, relay trên 2 Core hoạt động.
                    \item Triển khai app NOC \texttt{campus\_noc\_monitor.py}: Dashboard + REST API.
                \end{itemize}
            \end{block}
        \end{column}
        \begin{column}{0.5\textwidth}
            \begin{block}{Hạ tầng \& vận hành lab}
                \begin{itemize}
                    \item DHCP-Server (Win Server 2012 R2): 4 scope VLAN 10/20/30/40 Active.
                    \item OVS/Ryu tự khôi phục khi khởi động (systemd services).
                    \item Syslog (Kiwi) nhận log từ Core, Server Farm, FW-ASAv HA.
                    \item Quy trình đồng bộ một chiều host 1 $\rightarrow$ repo $\rightarrow$ host 2;
                          cấu hình nhúng trong file \texttt{.unl}.
                \end{itemize}
            \end{block}
        \end{column}
    \end{columns}
\end{frame}

% ============================================================
% UNFINISHED
% ============================================================
\section{Công việc chưa hoàn thành}

\begin{frame}{Công việc còn lại}
    \begin{columns}[T]
        \begin{column}{0.5\textwidth}
            \begin{alertblock}{Đang dở}
                \begin{itemize}
                    \item DHCP-Server nhận Discover nhưng \textbf{chưa trả OFFER}
                          $\rightarrow$ kiểm tra service, scope, giaddr, bắt gói 2 chiều.
                    \item Access-SW3/SW4 chưa kết nối ổn định với Ryu.
                    \item NOC Dashboard còn hiển thị tổng Rx/Tx = 0 do chưa có traffic qua datapath OVS.
                    \item Đưa bản patch app Ryu về repo \texttt{configs/}.
                \end{itemize}
            \end{alertblock}
        \end{column}
        \begin{column}{0.5\textwidth}
            \begin{block}{Kế hoạch tiếp theo}
                \begin{itemize}
                    \item Ping liên site 100 $\rightarrow$ 200/300/400 qua tunnel, kiểm tra route OMP/BGP.
                    \item Dịch vụ Web/Mail, truy cập đa site; hoàn thiện kịch bản ASDM.
                    \item Ma trận kiểm thử: DHCP từng site, failover FW-HA, cắt 1 màu WAN, mất controller SDN.
                    \item Hoàn thiện báo cáo, slide và kịch bản demo bảo vệ (dự kiến cuối 11/2026).
                \end{itemize}
            \end{block}
        \end{column}
    \end{columns}
\end{frame}

% ============================================================
% DIFFICULTIES
% ============================================================
\section{Khó khăn}

\begin{frame}{Khó khăn gặp phải}
    \begin{itemize}
        \item \textbf{Tài nguyên lab lớn:} hơn 50 node (IOL, ASAv, vEdge, Windows Server) trên EVE-NG
              $\Rightarrow$ khởi động chậm, node có thể không chạy sau khi start (vd.\ SDN\_CONTROLLER).
        \item \textbf{Trạng thái runtime bị mất sau reboot:} whitelist valid-vedges trên controller
              phải bổ sung lại, nếu không vEdge rơi khỏi fabric.
        \item \textbf{Lỗi khó chẩn đoán trên SD-WAN:} thiếu root-cert-chain, sai \texttt{color} trên
              tunnel-interface, route tới controller chỉ có 1 ngả $\Rightarrow$ tunnel down, health đỏ.
        \item \textbf{Hành vi OpenFlow:} table-miss gửi về controller làm rơi broadcast DHCP;
              Ryu bản cũ (Python 3.6) khác API so với tài liệu hiện hành.
        \item \textbf{Đồng bộ nhiều host:} cấu hình nhúng trong \texttt{.unl} dễ bị ghi đè khi
              export/import hoặc merge git $\Rightarrow$ phải chuẩn hoá quy trình một chiều.
        \item \textbf{Truy cập từ xa:} GUI vManage nằm trong lab, cần route và SSH tunnel để mở từ laptop.
    \end{itemize}
\end{frame}

\begin{frame}[plain]
    \centering
    \vfill
    {\Huge\color{primaryDark}\textbf{Cảm ơn thầy/cô đã lắng nghe!}}
    \vfill
\end{frame}

\end{document}
